\documentclass[11pt,a4paper]{article}

\usepackage[utf8]{inputenc}
\usepackage[T1]{fontenc}
\usepackage{lmodern}
\usepackage[a4paper,margin=2.5cm]{geometry}
\usepackage{setspace}
\usepackage{hyperref}
\usepackage{enumitem}
\usepackage{amsmath}

\setstretch{1.05}
\setlist[itemize]{nosep,leftmargin=2em}
\setlist[enumerate]{nosep,leftmargin=2em}

\title{Structured Rewrite Notes}
\author{}
\date{}

\begin{document}

\maketitle

\noindent\textbf{Source:} Kam \& Power, \textit{Principles of Physiology for the Anaesthetist}, 4th ed., Chapter 47-50: ``Basic Concepts of Acid--Base Physiology.''

\tableofcontents

\section{Chapter 47 --- Basic Concepts of Acid--Base Physiology}

\subsection{Core idea}

Acid--base physiology concerns the regulation of hydrogen ion concentration in body fluids. Despite continuous production of both volatile acid (from carbon dioxide) and non-volatile acids (from metabolism), plasma hydrogen ion concentration is normally maintained within a narrow range. This stability is essential because alterations in hydrogen ion concentration impair enzyme function, membrane excitability, cellular energy production, central nervous system function, and endocrine responses.

\subsection{Daily acid production and hydrogen ion balance}

\subsubsection{Volatile acid}

Aerobic metabolism generates a very large daily load of carbon dioxide. Carbon dioxide combines with water to form carbonic acid, which dissociates into hydrogen ions and bicarbonate. Under normal circumstances this does not create a net gain of hydrogen ions, because an equivalent amount of carbon dioxide is eliminated by the lungs.

\subsubsection{Non-volatile acid}

The body also produces fixed acids that cannot be excreted by the lungs. These include:

\begin{itemize}
    \item lactic acid
    \item ketone bodies
    \item sulphuric acid
    \item phosphoric acid
\end{itemize}

Although the daily quantity of fixed acid is much smaller than carbon dioxide production, it still constitutes an important acid load and must be handled largely by the kidneys.

\subsubsection{Gastrointestinal and renal contributions}

Hydrogen balance is also influenced by gastrointestinal losses and renal handling of acid and base.

\begin{itemize}
    \item Loss of acidic gastric fluid tends to remove hydrogen ions from the body.
    \item Loss of alkaline intestinal fluid tends to remove base.
    \item Urinary bicarbonate loss is equivalent to retaining hydrogen ions in the body.
    \item Urinary hydrogen ion excretion is equivalent to adding bicarbonate to the blood.
\end{itemize}

Normally, filtered bicarbonate is almost completely reabsorbed, and the kidneys excrete the daily fixed acid load.

\subsection{Definitions}

\subsubsection{Acid and base}

\begin{itemize}
    \item An \textbf{acid} is a proton donor.
    \item A \textbf{base} is a proton acceptor in solution.
\end{itemize}

\subsubsection{Hydrogen ion}

The hydrogen ion is effectively a proton. In aqueous solution it is present as the hydronium ion, \( \mathrm{H_3O^+} \).

\subsubsection{Dissociation and \(pK_a\)}

For a weak acid:
\[
\mathrm{HA \rightleftharpoons H^+ + A^-}
\]

The Henderson--Hasselbalch relationship is:
\[
\mathrm{pH = p}K_a + \log \left( \frac{[A^-]}{[HA]} \right)
\]

Important implications are:

\begin{itemize}
    \item the lower the \(pK_a\), the stronger the acid
    \item a buffer works best when the surrounding pH is close to its \(pK_a\)
\end{itemize}

\subsection{The pH system}

pH expresses hydrogen ion concentration on a logarithmic scale:
\[
\mathrm{pH = -\log_{10}[H^+]}
\]

Therefore:

\begin{itemize}
    \item a higher pH corresponds to a lower hydrogen ion concentration
    \item a lower pH corresponds to a higher hydrogen ion concentration
\end{itemize}

Because pH is logarithmic, equal changes in pH do not correspond to equal absolute changes in hydrogen ion concentration. A change of one pH unit represents a tenfold change in hydrogen ion concentration.

Body hydrogen ion concentration is normally maintained at approximately \(40\,\mathrm{nmol\,L^{-1}}\), corresponding to a pH close to 7.4.

\subsection{Buffers}

\subsubsection{Definition}

A buffer is a solution containing a weak acid and its conjugate base that resists a change in pH when acid or alkali is added.

\subsubsection{General principle}

A buffer binds added hydrogen ions reversibly:
\[
\mathrm{Buffer + H^+ \rightleftharpoons H.Buffer}
\]

Buffering is the first-line defence against acute pH change. It acts rapidly, but it does not remove acid from the body.

\subsubsection{Determinants of buffer effectiveness}

Buffer effectiveness depends on:

\begin{itemize}
    \item the quantity of buffer present
    \item its \(pK_a\)
    \item the pH of the surrounding solution
    \item whether the buffer operates in an open physiological system or a closed chemical system
\end{itemize}

\subsection{Acid--base homeostasis}

Plasma hydrogen ion concentration is maintained within a narrow range by three linked processes.

\subsubsection{Buffering}

Chemical buffer systems immediately minimize changes in hydrogen ion concentration.

\subsubsection{Compensation}

Physiological compensation restores the \( \mathrm{HCO_3^- / PCO_2} \) ratio towards normal. This is achieved mainly by the lungs and kidneys.

\subsubsection{Correction}

Correction refers to resolution of the primary disorder itself. This is the definitive step in restoring acid--base balance.

\subsection{The major buffer systems of the body}

The principal buffer systems are:

\begin{itemize}
    \item bicarbonate
    \item haemoglobin
    \item proteins
    \item phosphate
\end{itemize}

These may also be grouped as:

\begin{itemize}
    \item \textbf{bicarbonate buffer}
    \item \textbf{non-bicarbonate buffers:} haemoglobin, proteins, and phosphate
\end{itemize}

\subsection{Bicarbonate--carbonic acid buffer system}

\subsubsection{Why it matters}

The bicarbonate system is the most important extracellular buffer system and is central to clinical assessment of acid--base status.

\subsubsection{Components}

The system consists of:

\begin{itemize}
    \item carbonic acid (\(\mathrm{H_2CO_3}\)) as the weak acid
    \item bicarbonate (\(\mathrm{HCO_3^-}\)) as the conjugate base
\end{itemize}

Carbonic acid is generated from:
\[
\mathrm{CO_2 + H_2O \rightleftharpoons H_2CO_3}
\]

This reaction is accelerated by carbonic anhydrase, especially in red cells, renal tubules, and alveolar cells.

\subsubsection{Buffering of added acid}

When a strong acid is added:

\begin{itemize}
    \item the released hydrogen ions combine with bicarbonate
    \item carbonic acid is formed
    \item carbonic acid dissociates to carbon dioxide and water
    \item the lungs excrete the extra carbon dioxide
\end{itemize}

Thus bicarbonate provides the chemical buffering, while the lungs complete the physiological disposal of the acid load.

\subsubsection{Why this buffer is important despite its \(pK_a\)}

Its \(pK_a\) is relatively low compared with physiological pH, and the concentrations of its components are not exceptionally high. Nevertheless, the bicarbonate system is highly effective because:

\begin{itemize}
    \item \( \mathrm{PCO_2} \) can be rapidly regulated by ventilation
    \item \( \mathrm{HCO_3^-} \) can be regulated by the kidneys
\end{itemize}

It is therefore an open buffer system, which explains its dominance in extracellular fluid.

\subsubsection{Henderson and Henderson--Hasselbalch concepts}

The key clinical concept is that pH depends on the ratio of bicarbonate to dissolved carbon dioxide. The kidneys mainly regulate bicarbonate, while the lungs mainly regulate carbon dioxide.

\subsection{Haemoglobin as a buffer}

\subsubsection{Importance}

Haemoglobin is the principal non-bicarbonate buffer in blood and buffers both respiratory and metabolic acids.

\subsubsection{Why it is effective}

Haemoglobin contains many histidine residues with imidazole side chains whose \(pK_a\) is sufficiently close to physiological pH to make them useful buffers. In addition, haemoglobin is present at high concentration within red cells.

\subsubsection{Influence of oxygenation state}

\begin{itemize}
    \item Deoxyhaemoglobin is a better proton acceptor than oxyhaemoglobin.
    \item In tissue capillaries, oxygen unloading enhances buffering of hydrogen ions generated during carbon dioxide hydration.
\end{itemize}

This contributes to isohydric buffering and helps explain why venous blood is only slightly more acidic than arterial blood despite substantial carbon dioxide uptake.

\subsubsection{Carbamino formation}

Carbon dioxide can also bind directly to terminal amino groups on haemoglobin to form carbamino compounds. This accounts for part of the total carriage of carbon dioxide in blood.

\subsubsection{Clinical significance}

Although bicarbonate is the main plasma buffer, haemoglobin is the dominant buffer within whole blood, especially for carbonic acid.

\subsection{Proteins as buffers}

\subsubsection{General principle}

All proteins contain ionizable groups, but at physiological pH the main buffering contribution comes from histidine imidazole groups.

\subsubsection{Extracellular and intracellular roles}

\begin{itemize}
    \item In extracellular fluid, proteins contribute some buffering but less than bicarbonate or haemoglobin.
    \item In intracellular fluid, proteins are much more important because intracellular protein concentration is higher and intracellular pH is lower, closer to their effective buffering range.
\end{itemize}

\subsection{Phosphate buffer}

\subsubsection{Chemical basis}

Around physiological pH, the relevant phosphate pair is:

\begin{itemize}
    \item \( \mathrm{H_2PO_4^-} \)
    \item \( \mathrm{HPO_4^{2-}} \)
\end{itemize}

\subsubsection{Buffering value}

On chemical grounds, phosphate is potentially an effective buffer because its \(pK_a\) lies near physiological pH. However, its concentration in plasma is low, so its role in extracellular fluid is limited.

\subsubsection{Where phosphate matters}

Phosphate is more important:

\begin{itemize}
    \item intracellularly, where its concentration is higher
    \item in urine, where it contributes significantly to renal acid excretion
\end{itemize}

\subsubsection{Bone as alkali reserve}

During prolonged acidosis, calcium phosphate in bone becomes more soluble. Bone can therefore act as a longer-term alkali reserve.

\subsection{Distribution of buffers by body compartment}

\subsubsection{Blood}

The principal blood buffers are bicarbonate and haemoglobin.

Plasma bicarbonate is the main plasma buffer. Haemoglobin provides most of the red cell buffering capacity and most of the blood's ability to buffer carbonic acid. Plasma proteins and phosphate contribute relatively little.

\subsubsection{Interstitial fluid}

The main interstitial buffer is bicarbonate. Because interstitial fluid volume exceeds plasma volume, the interstitial compartment contributes substantially to buffering of non-carbonic acid.

\subsubsection{Intracellular fluid}

The major intracellular buffers are:

\begin{itemize}
    \item proteins
    \item phosphate
\end{itemize}

Their importance reflects both their concentration and the lower intracellular pH, which places them nearer their optimal buffering range.

\subsection{Whole-body (in vivo) versus blood-only (in vitro) titration}

\subsubsection{In vitro}

If carbon dioxide is added to blood in a closed container, bicarbonate generated during buffering remains within that blood sample.

\subsubsection{In vivo}

In the body, some bicarbonate produced in blood diffuses into the interstitial space. Therefore, after an acute rise in carbon dioxide:

\begin{itemize}
    \item plasma bicarbonate rises less
    \item arterial pH falls more
\end{itemize}

than would be predicted from an isolated blood sample alone.

\subsubsection{Why this matters}

The whole-body response reflects interaction between extracellular compartments. This distinction becomes most relevant at extremes of \( \mathrm{PCO_2} \) or in severe metabolic acidosis; in most routine clinical situations, the difference is small.

\section{Compensatory Mechanisms in Acid--Base Disorders}

\subsection{Fundamental principle}

Acid--base disturbances are limited by physiological responses that attempt to restore the ratio
\[
\frac{[HCO_3^-]}{PaCO_2}
\]
towards normal, thereby moving pH back in the normal direction.

Compensation is achieved by:
\begin{itemize}
  \item \textbf{respiratory adjustment} of CO\textsubscript{2} excretion
  \item \textbf{renal adjustment} of bicarbonate handling and acid excretion
\end{itemize}

Compensation is usually \textbf{incomplete}. It reduces the size of the pH disturbance, but does not normally restore pH completely to 7.4.

\subsection{Respiratory compensation}

Respiratory compensation alters \textbf{PaCO\textsubscript{2}} and is the main response to \textbf{primary metabolic disorders}.

\subsubsection{Metabolic acidosis}

\begin{itemize}
  \item Plasma bicarbonate falls.
  \item Blood hydrogen ion concentration rises.
  \item Ventilation is stimulated by medullary and carotid body chemoreceptors.
  \item \textbf{Hyperventilation lowers PaCO\textsubscript{2}}.
  \item This helps restore the \([HCO_3^-]/PaCO_2\) ratio.
\end{itemize}

Ventilation increases by about \textbf{twofold for a change in pH of 0.1 unit}.

\subsubsection{Metabolic alkalosis}

\begin{itemize}
  \item Plasma bicarbonate rises.
  \item Respiratory drive falls.
  \item \textbf{PaCO\textsubscript{2} rises} due to reduced ventilation.
  \item This tends to restore the ratio towards normal.
\end{itemize}

\subsubsection{Key features}

\begin{itemize}
  \item Respiratory compensation is \textbf{rapid}.
  \item It limits large acute changes in plasma hydrogen ion concentration.
  \item Its buffering capacity is about \textbf{twice that of the chemical buffers in extracellular fluid}.
\end{itemize}

\subsection{Renal compensation}

Renal compensation alters \textbf{plasma bicarbonate concentration} and is the main response to \textbf{primary respiratory disorders}.

\subsubsection{Time course and capacity}

\begin{itemize}
  \item It is \textbf{slow}.
  \item Maximal response takes about \textbf{7--10 days}.
  \item Maximum capacity is about \textbf{300 mmol H\textsuperscript{+}/day}.
\end{itemize}

\subsubsection{Daily acid load handled by the kidney}

\begin{itemize}
  \item About \textbf{70 mmol/day} of non-volatile acid is produced.
  \item About \textbf{4320 mmol/day} of bicarbonate is filtered.
  \item About \textbf{4320 mmol/day} of H\textsuperscript{+} must be secreted to reclaim this filtered bicarbonate.
  \item A further \textbf{70 mmol/day} of H\textsuperscript{+} is required to excrete the daily non-volatile acid load.
\end{itemize}

Total tubular H\textsuperscript{+} secretion is therefore about \textbf{4390 mmol/day}.

\subsubsection{Three renal components}

Renal regulation of H\textsuperscript{+} balance occurs by:
\begin{enumerate}
  \item \textbf{secretion of hydrogen ions with reabsorption of filtered bicarbonate}
  \item \textbf{excretion of titratable acid}
  \item \textbf{excretion of ammonia/ammonium}
\end{enumerate}

\subsection{Reabsorption of filtered bicarbonate}

This is the first major renal mechanism.

\subsubsection{Sites}

\begin{itemize}
  \item about \textbf{80\%} in the \textbf{proximal tubule}
  \item about \textbf{10\%} in the \textbf{thick ascending limb}
  \item the remainder in the \textbf{distal tubule and collecting duct}
\end{itemize}

\subsubsection{Mechanism}

Tubular cells generate H\textsuperscript{+} and HCO\textsubscript{3}\textsuperscript{-} from CO\textsubscript{2} and water via carbonic anhydrase.

Then:
\begin{itemize}
  \item H\textsuperscript{+} is secreted into the tubular lumen
  \item luminal H\textsuperscript{+} reacts with filtered bicarbonate
  \item CO\textsubscript{2} and water are formed
  \item CO\textsubscript{2} diffuses back into the cell
  \item intracellular carbonic anhydrase regenerates H\textsuperscript{+} and HCO\textsubscript{3}\textsuperscript{-}
  \item HCO\textsubscript{3}\textsuperscript{-} enters blood
  \item H\textsuperscript{+} is recycled into the lumen
\end{itemize}

\subsubsection{Functional meaning}

This process \textbf{reabsorbs filtered bicarbonate}, but it does \textbf{not create net new bicarbonate}. It is a reclamation process.

\subsection{Generation of new bicarbonate}

The kidney generates \textbf{new bicarbonate} when secreted H\textsuperscript{+} is excreted in urine buffered by \textbf{phosphate} or \textbf{ammonia/ammonium}, rather than merely being used to reclaim filtered bicarbonate.

This is essential because reabsorption of filtered bicarbonate alone cannot replace bicarbonate consumed buffering fixed acids.

\subsubsection{Phosphate mechanism: titratable acidity}

\paragraph{Definition}
Titratable acidity is the H\textsuperscript{+} bound to urinary buffers and equals the amount of alkali required to titrate urine back to \textbf{pH 7.4}.

\paragraph{Mechanism}
Once filtered bicarbonate has been fully reabsorbed:
\begin{itemize}
  \item additional secreted H\textsuperscript{+} combines with \textbf{monohydrogen phosphate}
  \item this forms \textbf{dihydrogen phosphate}
  \item for each dihydrogen phosphate formed, \textbf{one bicarbonate ion is added to the peritubular capillary blood}
\end{itemize}

That bicarbonate is \textbf{new bicarbonate}.

\paragraph{Features}
\begin{itemize}
  \item proximal tubule is the chief site
  \item other urinary buffers such as creatinine, $\beta$-hydroxybutyrate and sulphates contribute only a small amount
  \item phosphate buffering can excrete about \textbf{40 mmol/day} of H\textsuperscript{+}
  \item it is limited because phosphate availability cannot easily be increased
\end{itemize}

\subsubsection{Ammonium/ammonia mechanism}

This is the major adaptive system, especially in acidosis.

\paragraph{Why it matters}
The phosphate system alone is \textbf{not sufficient} to deal with normal daily H\textsuperscript{+} production, and cannot handle pathological acid loads adequately.

\paragraph{From glutamine to ammonium and bicarbonate}
In the proximal tubule, thick ascending limb and distal tubules:
\begin{itemize}
  \item glutamine is metabolized
  \item this produces \textbf{2 NH\textsubscript{4}\textsuperscript{+}} and $\alpha$-ketoglutarate
  \item $\alpha$-ketoglutarate is converted to \textbf{2 bicarbonate ions}
  \item NH\textsubscript{4}\textsuperscript{+} is secreted into the lumen
  \item HCO\textsubscript{3}\textsuperscript{-} is reabsorbed into blood
\end{itemize}

This means ammonium production is linked directly to formation of \textbf{new bicarbonate}.

\paragraph{Ammonia trapping in the collecting duct}
\begin{itemize}
  \item NH\textsubscript{3} diffuses into the collecting duct lumen
  \item secreted H\textsuperscript{+} combines with NH\textsubscript{3} to form \textbf{NH\textsubscript{4}\textsuperscript{+}}
  \item NH\textsubscript{4}\textsuperscript{+} cannot diffuse back and is excreted
\end{itemize}

Thus NH\textsubscript{4}\textsuperscript{+} acts as a \textbf{sink for H\textsuperscript{+}}.

\paragraph{Quantitative importance}
\begin{itemize}
  \item normal NH\textsubscript{4}\textsuperscript{+} excretion: about \textbf{30--50 mmol/day}
  \item severe metabolic acidosis: may rise to \textbf{>300 mmol/day}
  \item this is accompanied by an equivalent increase in \textbf{new bicarbonate} production/reabsorption
\end{itemize}

\subsection{Renal response in the major primary disorders}

\subsubsection{Respiratory acidosis}

Raised PaCO\textsubscript{2} and raised extracellular H\textsuperscript{+} stimulate:
\begin{itemize}
  \item renal H\textsuperscript{+} secretion
  \item ammonium secretion
  \item bicarbonate reabsorption
\end{itemize}

Result:
\begin{itemize}
  \item plasma bicarbonate rises
  \item pH is partially corrected
\end{itemize}

\subsubsection{Respiratory alkalosis}

Low PaCO\textsubscript{2} inhibits:
\begin{itemize}
  \item H\textsuperscript{+} secretion
  \item ammonium secretion
\end{itemize}

Result:
\begin{itemize}
  \item less bicarbonate is reabsorbed
  \item plasma bicarbonate falls
  \item pH moves back towards normal
\end{itemize}

\subsubsection{Metabolic acidosis}

The main immediate compensation is \textbf{hyperventilation}, but renal acid excretion and new bicarbonate generation also become important in sustained states.

\subsubsection{Metabolic alkalosis}

The main compensation is \textbf{hypoventilation}, while renal handling shifts towards bicarbonate loss if conditions allow.

\subsection{Clinical effects of acid--base changes}

Acid--base disturbances have widespread \textbf{physiological and biochemical effects}, either directly on organs or indirectly via the \textbf{autonomic} and \textbf{endocrine} systems.

\subsubsection{Cardiovascular effects}

\paragraph{Acidosis}
\begin{itemize}
  \item direct \textbf{negative inotropic} effect on myocardium
  \item due to inhibition of the slow inward calcium current and reduced calcium release from the sarcoplasmic reticulum
  \item \textbf{acute respiratory acidosis} depresses contractility more than acute metabolic acidosis
  \item acidosis increases catecholamine release, initially producing \textbf{tachycardia}
  \item if \textbf{pH falls below 7.2}, myocardial depression predominates
  \item atrial and ventricular arrhythmias are more common
  \item acidosis directly causes vasodilatation in \textbf{skin, skeletal muscle, uterus and heart}
  \item but causes \textbf{pulmonary vasoconstriction}
  \item with mild acidosis (\textbf{pH 7.2--7.3}), raised catecholamines may produce \textbf{systemic, renal and splanchnic vasoconstriction}
\end{itemize}

\paragraph{Alkalosis}
\begin{itemize}
  \item increases \textbf{coronary and systemic vascular resistance}
  \item shifts the oxyhaemoglobin dissociation curve \textbf{to the left}
  \item therefore impairs oxygen delivery to tissues
\end{itemize}

\subsubsection{Respiratory effects}

\paragraph{Acidosis}
\begin{itemize}
  \item increases minute ventilation in metabolic acidosis through central chemoreceptor stimulation
  \item in respiratory acidosis, minute ventilation rises by \textbf{2--3 L/min for every 1 mmHg rise in PaCO\textsubscript{2}}
  \item the response to respiratory acidosis is faster than to metabolic acidosis because CO\textsubscript{2} crosses the blood--brain barrier more readily than H\textsuperscript{+}
  \item hyperventilation peaks at a \textbf{PaCO\textsubscript{2} of 100 mmHg (13.3 kPa)}
  \item beyond this, further rises in PaCO\textsubscript{2} cause \textbf{respiratory depression}
  \item hypercapnia causes \textbf{direct bronchodilatation}, though vagal effects may cause bronchoconstriction
\end{itemize}

\paragraph{Oxyhaemoglobin dissociation curve}
\begin{itemize}
  \item acidosis shifts it \textbf{right} and improves oxygen unloading
  \item alkalosis shifts it \textbf{left} and reduces oxygen delivery
\end{itemize}

\subsubsection{Electrolyte effects}

\paragraph{Calcium}
\begin{itemize}
  \item in acidosis, H\textsuperscript{+} competes with calcium for albumin binding sites
  \item ionized calcium therefore \textbf{rises}
  \item the reverse occurs in alkalosis, so ionized calcium \textbf{falls}
\end{itemize}

\paragraph{Potassium}
\begin{itemize}
  \item in acidosis, H\textsuperscript{+} moves into cells and K\textsuperscript{+} moves out into extracellular fluid
  \item estimated change: about \textbf{0.6 mmol/L} in serum K\textsuperscript{+} for each \textbf{0.1 unit} change in pH
\end{itemize}

\paragraph{Bone}
\begin{itemize}
  \item chronic metabolic acidosis can mobilize calcium from bone and raise serum calcium
\end{itemize}

\subsubsection{Central nervous system}

\begin{itemize}
  \item severe acidosis can impair consciousness
  \item CNS effects are mediated largely by changes in \textbf{cerebral blood flow} and \textbf{intracranial pressure}
  \item alkalosis may precipitate \textbf{epilepsy}
\end{itemize}

\subsubsection{Gastrointestinal effects}

\begin{itemize}
  \item acidosis may cause \textbf{splanchnic vasoconstriction}
  \item gastrointestinal motility may be \textbf{reduced}
\end{itemize}

\subsection{High-yield distinction: bicarbonate reabsorption versus new bicarbonate generation}

\subsubsection{Reabsorption of filtered bicarbonate}
\begin{itemize}
  \item prevents bicarbonate loss
  \item occurs mainly in proximal tubule
  \item requires H\textsuperscript{+} secretion
  \item \textbf{does not produce net new bicarbonate}
\end{itemize}

\subsubsection{Generation of new bicarbonate}
\begin{itemize}
  \item occurs when H\textsuperscript{+} is excreted as:
  \begin{itemize}
    \item \textbf{titratable acid} (mainly phosphate)
    \item \textbf{NH\textsubscript{4}\textsuperscript{+}}
  \end{itemize}
  \item each H\textsuperscript{+} excreted in these forms adds bicarbonate to blood
  \item this replaces bicarbonate consumed buffering non-volatile acids
\end{itemize}

\subsection{Exam summary}

\begin{itemize}
  \item Compensation aims to restore the \([HCO_3^-]/PaCO_2\) ratio and move pH towards normal.
  \item \textbf{Metabolic disorders} are compensated mainly by the \textbf{lungs}.
  \item \textbf{Respiratory disorders} are compensated mainly by the \textbf{kidneys}.
  \item Respiratory compensation is \textbf{rapid}; renal compensation is \textbf{slow}.
  \item Renal compensation involves:
  \begin{enumerate}
    \item bicarbonate reclamation
    \item titratable acid excretion
    \item ammonium excretion
  \end{enumerate}
  \item \textbf{New bicarbonate} is generated only when H\textsuperscript{+} is excreted as phosphate or ammonium.
  \item Acid--base disorders have major \textbf{cardiovascular, respiratory, electrolyte, neurological and gastrointestinal} effects.
\end{itemize}

\section{Clinical Aspects of Acid--Base Control}

\subsection{Core terminology}

\subsubsection{Acidaemia}
Arterial blood pH below the normal range, defined as \(\mathrm{pH} < 7.36\), corresponding to \([\mathrm{H}^+] > 44~\mathrm{nmol/L}\).

\subsubsection{Alkalaemia}
Arterial blood pH above the normal range, defined as \(\mathrm{pH} > 7.44\), corresponding to \([\mathrm{H}^+] < 36~\mathrm{nmol/L}\).

\subsubsection{Acidosis}
A pathological process that tends to lower arterial pH. It is classified as:
\begin{itemize}
  \item \textbf{Respiratory}, when the primary abnormality is in \(\mathrm{PaCO_2}\)
  \item \textbf{Metabolic}, when the primary abnormality is a reduction in bicarbonate due to fixed/non-volatile acid gain or bicarbonate loss
\end{itemize}

\subsubsection{Alkalosis}
A pathological process that tends to raise arterial pH.

\subsubsection{Compensation}
A physiological response, either respiratory or renal, that attempts to return arterial pH toward normal. Compensation reduces the pH disturbance but usually does not completely normalize pH.

\subsection{How acid--base disorders are assessed}

The essential measurements are:
\begin{itemize}
  \item pH
  \item \(\mathrm{PaCO_2}\)
  \item Calculated plasma \(\mathrm{HCO_3^-}\)
\end{itemize}

These allow identification of:
\begin{itemize}
  \item the primary disorder
  \item the degree of compensation
\end{itemize}

\subsection{Historical and laboratory approaches}

\subsubsection{Older ``alkali reserve'' approach}
Historically, plasma was titrated to a \(\mathrm{PaCO_2}\) of \(40~\mathrm{mmHg}\) (\(5.3~\mathrm{kPa}\)) and total \(\mathrm{CO_2}\) content measured. Bicarbonate was then measured and pH derived using the Henderson--Hasselbalch equation.

\textbf{Limitation:} this method ignored the buffering contribution of haemoglobin, so it was incomplete.

\subsection{In vitro assessment methods}

\subsubsection{Astrup method}
This is an \emph{in vitro} equilibration method used to determine \(\mathrm{PaCO_2}\).

Process:
\begin{itemize}
  \item initial blood pH is measured
  \item two blood samples are equilibrated with \(4\%\) and \(8\%\) \(\mathrm{CO_2}\) in oxygen at \(37^\circ\mathrm{C}\)
  \item pH is measured at each known \(\mathrm{CO_2}\) tension
  \item the values are plotted against \(\log \mathrm{PaCO_2}\)
  \item the original sample's \(\mathrm{PaCO_2}\) is derived from the graph
\end{itemize}

This method defines the pH--\(\mathrm{PaCO_2}\) relationship for that blood sample.

\subsubsection{Standard bicarbonate}
Standard bicarbonate is the bicarbonate concentration measured in fully oxygenated whole blood at:
\begin{itemize}
  \item \(\mathrm{PaCO_2} = 40~\mathrm{mmHg}\) (\(5.3~\mathrm{kPa}\))
  \item \(37^\circ\mathrm{C}\)
\end{itemize}

Normal range:
\begin{itemize}
  \item \(22\text{--}26~\mathrm{mmol/L}\)
\end{itemize}

\textbf{Limitation:} it does not account for buffering in the interstitial and intracellular compartments.

\subsubsection{Base excess}
Base excess is based on the idea that when blood pH is normal, the ratios and total concentrations of non-carbonic buffers are normal.

It estimates the amount of acid or base that must be added to return blood to:
\begin{itemize}
  \item \(\mathrm{pH} = 7.4\)
  \item under standard conditions
\end{itemize}

Conceptually:
\begin{itemize}
  \item the non-carbonic buffers are restored to normal
  \item the remaining bicarbonate change reflects the metabolic acid or alkali load
\end{itemize}

\textbf{Advantage:} simple and clinically useful.

\textbf{Main limitation:} it is still an \emph{in vitro} system and does not include extravascular buffering.

\subsubsection{Buffer base}
Buffer base is the total concentration of all buffer anions in blood, including:
\begin{itemize}
  \item bicarbonate
  \item haemoglobin
  \item protein
  \item phosphate
\end{itemize}

This is another attempt to quantify the metabolic component of acid--base status.

\subsection{Limitation of in vitro methods}

\emph{In vitro} methods assess primarily the blood compartment, which represents only about one-third of total body buffering capacity.

They do not fully account for movement of:
\begin{itemize}
  \item \(\mathrm{H^+}\)
  \item \(\mathrm{CO_2}\)
  \item \(\mathrm{HCO_3^-}\)
\end{itemize}

between:
\begin{itemize}
  \item blood
  \item interstitial fluid
  \item intracellular fluid
\end{itemize}

As a result, \emph{in vitro} approaches tend to overestimate the magnitude of whole-body acid--base change.

\subsection{In vivo evaluation}

\subsubsection{Why in vivo assessment matters}
\emph{In vivo} titration curves are derived from actual acute and chronic acid--base disorders. Their main advantage is that they reflect:
\begin{itemize}
  \item interaction between body compartments
  \item real physiological compensation
\end{itemize}

This makes them more clinically representative than isolated blood-based systems.

\subsection{In vivo patterns in respiratory disorders}

\subsubsection{Acute respiratory acidosis}
A sudden rise in \(\mathrm{PaCO_2}\) increases carbonic acid and raises both \([\mathrm{H^+}]\) and \([\mathrm{HCO_3^-}]\).

For each \(10~\mathrm{mmHg}\) (\(1.3~\mathrm{kPa}\)) increase in \(\mathrm{PaCO_2}\) above \(40~\mathrm{mmHg}\):
\begin{itemize}
  \item \(\mathrm{HCO_3^-}\) rises by \(0.08~\mathrm{mmol/L}\)
  \item pH falls by \(0.07\)
  \item equivalent to about \(+8~\mathrm{nmol/L}\) \([\mathrm{H^+}]\)
\end{itemize}

\subsubsection{Chronic respiratory acidosis}
With time, renal compensation increases \(\mathrm{H^+}\) excretion and bicarbonate reabsorption.

For each \(10~\mathrm{mmHg}\) (\(1.3~\mathrm{kPa}\)) rise in \(\mathrm{PaCO_2}\) above \(40~\mathrm{mmHg}\):
\begin{itemize}
  \item \(\mathrm{HCO_3^-}\) rises by \(4~\mathrm{mmol/L}\)
  \item pH falls by \(0.03\)
  \item equivalent to about \(+3.2~\mathrm{nmol/L}\) \([\mathrm{H^+}]\)
\end{itemize}

\subsubsection{Acute respiratory alkalosis}
A sudden fall in \(\mathrm{PaCO_2}\) reduces \([\mathrm{H^+}]\) and \([\mathrm{HCO_3^-}]\).

For each \(10~\mathrm{mmHg}\) (\(1.3~\mathrm{kPa}\)) fall in \(\mathrm{PaCO_2}\) below \(40~\mathrm{mmHg}\):
\begin{itemize}
  \item \(\mathrm{HCO_3^-}\) falls by \(2~\mathrm{mmol/L}\)
  \item pH rises by \(0.08\)
  \item equivalent to about \(-8~\mathrm{nmol/L}\) \([\mathrm{H^+}]\)
\end{itemize}

\subsubsection{Chronic respiratory alkalosis}
Renal compensation reduces \(\mathrm{H^+}\) secretion and bicarbonate reabsorption.

For each \(10~\mathrm{mmHg}\) (\(1.3~\mathrm{kPa}\)) fall in \(\mathrm{PaCO_2}\) below \(40~\mathrm{mmHg}\):
\begin{itemize}
  \item \(\mathrm{HCO_3^-}\) falls by \(6~\mathrm{mmol/L}\)
  \item pH rises by \(0.03\)
  \item equivalent to about \(-3.2~\mathrm{nmol/L}\) \([\mathrm{H^+}]\)
\end{itemize}

\subsection{In vivo patterns in metabolic disorders}

\subsubsection{Metabolic acidosis}
Ventilation rises rapidly, but respiratory compensation is only partial.

\emph{In vivo} titration data indicate:
\begin{itemize}
  \item \(\mathrm{PaCO_2}\) falls by \(0.1~\mathrm{mmHg}\) (\(0.01~\mathrm{kPa}\)) for each \(1~\mathrm{mmol/L}\) fall in \(\mathrm{HCO_3^-}\)
  \item pH falls by \(0.02\)
\end{itemize}

\subsubsection{Metabolic alkalosis}
Respiratory drive falls, so \(\mathrm{PaCO_2}\) rises, partially restoring the \(\mathrm{HCO_3^-}/\mathrm{PaCO_2}\) ratio.

Predicted \(\mathrm{PaCO_2}\):
\[
\mathrm{PaCO_2} = 0.7 \times [\mathrm{HCO_3^-}] + 20 \pm 2~\mathrm{mmHg}
\]

\subsection{Temperature and acid--base control}

\subsubsection{Effect of cooling}
The dissociation of water is temperature dependent. As temperature falls:
\begin{itemize}
  \item the pH of neutral water rises by \(0.017\) per \(1^\circ\mathrm{C}\)
  \item blood pH rises by about \(0.015\) per \(1^\circ\mathrm{C}\)
  \item \(\mathrm{PaCO_2}\) falls by about \(4.5\%\)
\end{itemize}

Thus hypothermia produces an apparent alkalinizing effect.

\subsection{Acid--base strategies during hypothermia}

\subsubsection{pH-stat management}
In pH-stat management, blood gas values are corrected so that during cooling:
\begin{itemize}
  \item pH is maintained at \(7.4\)
  \item \(\mathrm{PaCO_2}\) is maintained at \(40~\mathrm{mmHg}\)
\end{itemize}

This means \(\mathrm{CO_2}\) is effectively added during hypothermia to keep measured values ``normal'' at the patient's actual temperature.

Implications described in the source:
\begin{itemize}
  \item produces relative acidosis compared with \(\alpha\)-stat
  \item higher \(\mathrm{PaCO_2}\) increases cerebral blood flow via vasodilatation
  \item may impair autoregulation
  \item may increase risk of cerebral embolism and postoperative cerebral dysfunction
\end{itemize}

\subsubsection{\(\alpha\)-stat management}
With \(\alpha\)-stat management:
\begin{itemize}
  \item blood gases are measured at \(37^\circ\mathrm{C}\)
  \item values are not corrected for the patient's lower body temperature
\end{itemize}

Therefore, during hypothermia:
\begin{itemize}
  \item blood pH is allowed to rise
  \item the alkalosis is left uncorrected
\end{itemize}

This is based on the idea that the ionization state of histidine imidazole groups remains appropriate across temperature changes, preserving enzyme function and cerebral autoregulation.

Claimed advantages in the source:
\begin{itemize}
  \item better preservation of enzyme function
  \item more normal cerebral autoregulation
  \item improved neurological outcomes
\end{itemize}

Differences between \(\alpha\)-stat and pH-stat become especially important below \(25^\circ\mathrm{C}\), and are less relevant above \(32^\circ\mathrm{C}\).

\subsection{Biochemical description of an acid--base defect}

Any acid--base abnormality can be described using three parameters:
\begin{itemize}
  \item pH, showing overall acidity/alkalinity
  \item \(\mathrm{PaCO_2}\), representing the respiratory component
  \item \(\mathrm{HCO_3^-}\), representing the metabolic component
\end{itemize}

pH and \(\mathrm{PaCO_2}\) are measured directly. \(\mathrm{HCO_3^-}\) is calculated from the Henderson--Hasselbalch equation.

Derived values such as:
\begin{itemize}
  \item standard bicarbonate
  \item base excess
\end{itemize}

may help, but they do not fully reflect \emph{in vivo} physiology.

For most clinical purposes, assessment can be based on:
\begin{itemize}
  \item pH
  \item \(\mathrm{PaCO_2}\)
  \item \(\mathrm{HCO_3^-}\)
\end{itemize}

\subsection{Graphical evaluation}

\subsubsection{pH--\(\mathrm{PaCO_2}\) diagram}
A pH/\(\mathrm{PaCO_2}\) diagram demonstrates the \emph{in vivo} relationship between pH and \(\mathrm{PaCO_2}\) in primary acid--base disorders.

It allows recognition of:
\begin{itemize}
  \item normal range
  \item acute respiratory acidosis
  \item acute respiratory alkalosis
  \item chronic respiratory compensation
\end{itemize}

This helps separate acute from chronic respiratory disorders.

\subsubsection{pH--\(\mathrm{HCO_3^-}\) diagram}
A pH/\(\mathrm{HCO_3^-}\) diagram classifies disorders according to position relative to:
\begin{itemize}
  \item the \(\mathrm{PaCO_2}\) \(40~\mathrm{mmHg}\) isobar
  \item the buffer line
\end{itemize}

Interpretation:
\begin{itemize}
  \item respiratory disorders lie to either side of the \(\mathrm{PaCO_2}\) \(40~\mathrm{mmHg}\) isobar
  \item metabolic disorders lie above or below the buffer line
\end{itemize}

Also:
\begin{itemize}
  \item standard bicarbonate \(< 24~\mathrm{mmol/L}\) suggests metabolic acidosis
  \item standard bicarbonate \(> 24~\mathrm{mmol/L}\) suggests metabolic alkalosis
\end{itemize}

% Fragment for \input{} / \include{}

\section{Anion Gap and Stewart's Strong Ion Difference}

\subsection{Overview}

Two complementary frameworks are commonly used to analyse metabolic acid--base disorders:
\begin{enumerate}[label=\arabic*.]
    \item the \textbf{anion gap} approach, which helps detect the presence of unmeasured anions;
    \item the \textbf{Stewart physicochemical} approach, which explains acid--base disturbances using independent physicochemical variables.
\end{enumerate}

The anion gap is clinically useful as a rapid bedside tool, whereas Stewart's model provides a more mechanistic explanation of how electrolyte and protein changes alter pH.

\subsection{The Anion Gap}

\subsubsection{Definition}

The anion gap estimates the difference between commonly measured plasma cations and anions:
\[
\text{Anion gap} = \bigl[\mathrm{Na}^+\bigr] + \bigl[\mathrm{K}^+\bigr] - \bigl(\bigl[\mathrm{Cl}^-\bigr] + \bigl[\mathrm{HCO}_3^-\bigr]\bigr)
\]

A typical normal range is approximately \(5\text{--}15~\mathrm{mEq/L}\).

\subsubsection{Physiological basis}

Plasma must remain electrically neutral. The measured cations do not exactly equal the measured anions because some ions are not routinely included in the calculation. The anion gap therefore reflects the contribution of \textbf{unmeasured anions}, chiefly:
\begin{itemize}
    \item albumin;
    \item phosphate.
\end{itemize}

\subsubsection{Interpretation}

A metabolic acidosis may be divided into:
\begin{itemize}
    \item \textbf{high anion gap metabolic acidosis}, where bicarbonate falls and is replaced by unmeasured anions;
    \item \textbf{normal anion gap (hyperchloraemic) metabolic acidosis}, where bicarbonate falls but is replaced mainly by chloride.
\end{itemize}

\subsubsection{High anion gap metabolic acidosis}

A raised anion gap suggests accumulation of acids whose anions are not routinely measured. Important causes include:
\begin{itemize}
    \item lactic acidosis;
    \item ketoacidosis;
    \item uraemia or renal failure;
    \item salicylate poisoning;
    \item methanol poisoning;
    \item ethylene glycol poisoning.
\end{itemize}

Examples of associated anions include lactate, acetoacetate, \(\beta\)-hydroxybutyrate, and oxalate.

\subsubsection{Normal anion gap metabolic acidosis}

This usually reflects bicarbonate loss or chloride gain. Causes include:
\begin{itemize}
    \item diarrhoea;
    \item pancreatic fistula;
    \item renal tubular acidosis;
    \item hydrochloric acid administration;
    \item ammonium chloride administration;
    \item acetazolamide.
\end{itemize}

\subsubsection{Low anion gap}

A low anion gap may occur in states of reduced plasma protein concentration, especially hypoalbuminaemia.

\subsubsection{Correction for albumin}

Because albumin contributes substantially to the normal anion gap, hypoalbuminaemia can mask the presence of unmeasured anions. A correction can therefore be applied:
\[
\text{Corrected anion gap} = \text{Anion gap} + 0.25 \times (40 - \text{albumin})
\]

where albumin is expressed in \(\mathrm{g/L}\), assuming a normal value of \(40~\mathrm{g/L}\).

\subsection{Limitations of the Anion Gap Approach}

The anion gap is simple and clinically valuable, but it has limitations:
\begin{itemize}
    \item it is influenced by albumin concentration;
    \item it does not provide a full physicochemical explanation for acid--base disorders;
    \item it may miss mixed disturbances unless interpreted in context.
\end{itemize}

\subsection{Stewart's Physicochemical Approach}

\subsubsection{Rationale}

Stewart proposed that acid--base chemistry should be understood from first principles rather than by treating bicarbonate as an independent metabolic variable. In this model, hydrogen ion concentration is determined by the behaviour of water and by laws governing electrical neutrality and chemical equilibrium.

Water dissociates as:
\[
\mathrm{H_2O} \rightleftharpoons \mathrm{H}^+ + \mathrm{OH}^-
\]

A solution is:
\begin{itemize}
    \item \textbf{neutral} when \([\mathrm{H}^+] = [\mathrm{OH}^-]\),
    \item \textbf{acidic} when \([\mathrm{H}^+] > [\mathrm{OH}^-]\),
    \item \textbf{alkaline} when \([\mathrm{H}^+] < [\mathrm{OH}^-]\).
\end{itemize}

\subsubsection{Strong and weak ions}

\paragraph{Strong ions.}
Strong ions dissociate almost completely in solution. In plasma these include:
\begin{itemize}
    \item \(\mathrm{Na}^+\),
    \item \(\mathrm{K}^+\),
    \item \(\mathrm{Ca}^{2+}\),
    \item \(\mathrm{Mg}^{2+}\),
    \item \(\mathrm{Cl}^-\),
    \item other strong anions such as lactate.
\end{itemize}

\paragraph{Weak acids.}
Weak acids are only partially dissociated. The principal plasma weak acids are:
\begin{itemize}
    \item albumin;
    \item phosphate.
\end{itemize}

\subsubsection{The three independent variables}

According to Stewart, the acid--base state is determined by three independent variables:
\begin{enumerate}[label=\arabic*.]
    \item \textbf{\(P_{\mathrm{CO_2}}\)} -- the respiratory component;
    \item \textbf{strong ion difference (SID)} -- the difference between fully dissociated cations and anions;
    \item \textbf{total weak acids (\(A_{\mathrm{tot}}\))} -- mainly albumin and phosphate.
\end{enumerate}

\subsection{Strong Ion Difference}

\subsubsection{Definition}

The apparent strong ion difference is:
\[
\mathrm{SID_a} =
\bigl[\mathrm{Na}^+\bigr] + \bigl[\mathrm{K}^+\bigr] + \bigl[\mathrm{Ca}^{2+}\bigr] + \bigl[\mathrm{Mg}^{2+}\bigr]
-
\bigl(\bigl[\mathrm{Cl}^-\bigr] + \text{other strong anions}\bigr)
\]

A typical normal range is approximately \(40\text{--}44~\mathrm{mEq/L}\).

\subsubsection{Interpretation}

\begin{itemize}
    \item A \textbf{decreased SID} promotes metabolic acidosis.
    \item An \textbf{increased SID} promotes metabolic alkalosis.
\end{itemize}

The key principle is that changes in the balance between strong cations and strong anions alter water dissociation and therefore alter hydrogen ion concentration.

\subsubsection{Causes of a decreased SID}

A lower SID occurs when strong anions rise or strong cations fall. Examples include:
\begin{itemize}
    \item hyperchloraemia after large-volume \(0.9\%\) saline infusion;
    \item hyponatraemia;
    \item accumulation of lactate.
\end{itemize}

This produces a metabolic acidosis.

\subsubsection{Causes of an increased SID}

A higher SID occurs when strong cations rise or strong anions fall. Examples include:
\begin{itemize}
    \item hypochloraemia due to vomiting;
    \item hypochloraemia due to loop diuretics;
    \item hypernatraemia following sodium bicarbonate administration.
\end{itemize}

This produces a metabolic alkalosis.

\subsection{Total Weak Acids (\(A_{\mathrm{tot}}\))}

\subsubsection{Definition}

\(A_{\mathrm{tot}}\) represents the total concentration of plasma weak acids, mainly:
\begin{itemize}
    \item albumin;
    \item phosphate.
\end{itemize}

A typical normal range is approximately \(12\text{--}24~\mathrm{mEq/L}\).

\subsubsection{Interpretation}

\begin{itemize}
    \item \textbf{Increased \(A_{\mathrm{tot}}\)} causes metabolic acidosis.
    \item \textbf{Decreased \(A_{\mathrm{tot}}\)} causes metabolic alkalosis.
\end{itemize}

\subsubsection{Clinical examples}

\begin{itemize}
    \item Hyperphosphataemia in chronic renal failure increases \(A_{\mathrm{tot}}\) and tends to produce acidosis.
    \item Hypoalbuminaemia, common in long-stay ICU patients, decreases \(A_{\mathrm{tot}}\) and tends to produce alkalosis.
\end{itemize}

\subsection{Physicochemical Principles Underpinning Stewart's Model}

Stewart's analysis rests on three principles:
\begin{enumerate}[label=\arabic*.]
    \item \textbf{Electroneutrality}: the sum of positive charges must equal the sum of negative charges.
    \item \textbf{Mass action}: dissociation equilibria must be satisfied.
    \item \textbf{Conservation of mass}: the total amount of each substance equals the sum of its dissociated and undissociated forms.
\end{enumerate}

\subsection{Strong Ion Gap}

\subsubsection{Concept}

The strong ion gap (SIG) estimates the presence of unmeasured ions within the Stewart framework.

The \textbf{effective strong ion difference} (\(\mathrm{SID_e}\)) is derived from bicarbonate and weak acids. Under normal circumstances:
\[
\mathrm{SID_a} = \mathrm{SID_e}
\]

If unmeasured anions are present:
\[
\mathrm{SID_a} > \mathrm{SID_e}
\]

The difference between them is the \textbf{strong ion gap}.

\subsubsection{Clinical significance}

A raised SIG suggests additional unmeasured anions, for example:
\begin{itemize}
    \item lactate;
    \item acetoacetate;
    \item \(\beta\)-hydroxybutyrate;
    \item other pathological organic anions.
\end{itemize}

\subsection{Organ Contributions to Acid--Base Balance}

\subsubsection{Lungs}

The lungs regulate \(P_{\mathrm{CO_2}}\), thereby determining the respiratory component of acid--base balance.

\subsubsection{Kidneys}

The kidneys alter electrolyte composition, especially chloride handling, and therefore influence SID.

\subsubsection{Liver and gut}

The liver and gastrointestinal tract can influence the concentration of weak acids and therefore modify \(A_{\mathrm{tot}}\).

\subsection{Clinical Applications}

\subsubsection{Vomiting}

Vomiting causes loss of gastric hydrochloric acid. Chloride falls relative to sodium, SID increases, and metabolic alkalosis develops.

\subsubsection{Large-volume saline infusion}

\(0.9\%\) saline contains a chloride concentration higher than that of plasma. Administration of large volumes raises plasma chloride, lowers SID, and causes a hyperchloraemic metabolic acidosis.

\subsubsection{Sodium bicarbonate administration}

In Stewart terms, the alkalinising effect is explained primarily by the sodium load increasing SID rather than by bicarbonate acting as an independent variable.

\subsubsection{Hypoalbuminaemia}

Reduced albumin lowers \(A_{\mathrm{tot}}\) and therefore tends to produce metabolic alkalosis.

\subsection{Anion Gap versus Stewart: Comparison}

\subsubsection{Anion gap approach}

\begin{itemize}
    \item Simple and rapid.
    \item Useful for detecting unmeasured anions.
    \item Particularly valuable in metabolic acidosis.
    \item Must be corrected for albumin when appropriate.
\end{itemize}

\subsubsection{Stewart approach}

\begin{itemize}
    \item More mechanistic and comprehensive.
    \item Explains acid--base disorders in terms of \(P_{\mathrm{CO_2}}\), SID, and \(A_{\mathrm{tot}}\).
    \item Particularly helpful for understanding saline-induced hyperchloraemic acidosis, chloride-depletion alkalosis, and hypoalbuminaemic alkalosis.
\end{itemize}

\subsection{Key Revision Points}

\begin{enumerate}[label=\arabic*.]
    \item The anion gap estimates the presence of unmeasured anions.
    \item A high anion gap metabolic acidosis suggests lactate, ketones, uraemia, or toxins.
    \item A normal anion gap metabolic acidosis usually reflects bicarbonate loss or chloride gain.
    \item Stewart's three independent variables are \(P_{\mathrm{CO_2}}\), SID, and \(A_{\mathrm{tot}}\).
    \item A low SID causes acidosis; a high SID causes alkalosis.
    \item A high \(A_{\mathrm{tot}}\) causes acidosis; a low \(A_{\mathrm{tot}}\) causes alkalosis.
    \item A raised SIG indicates additional unmeasured anions.
\end{enumerate}

\end{document}
